% Figure: the literature catalog. Data: data/lit-top.csv and data/lit-best.csv, both written by
% data/regen_literature.py from docs/assets/literature/literature.db.
\pgfplotstableread[col sep=comma]{data/lit-top.csv}\littop
\pgfplotstableread[col sep=comma]{data/lit-best.csv}\litbest
\begin{tikzpicture}
\begin{axis}[
    ttaxis,
    name=top,
    width=0.47\linewidth, height=6.4cm,
    xbar, bar width=6pt,
    xmin=0,
    xlabel={papers using the dataset},
    ytick=data, y dir=reverse,
    yticklabels from table={\littop}{dataset},
    yticklabel style={font=\ttfamily\scriptsize, color=ttink},
    ymajorgrids=false,
    enlarge y limits=0.06,
    nodes near coords,
    nodes near coords style={font=\sffamily\tiny, color=ttdeep},
    title={Most-used datasets},
  ]
  \addplot[ttbar] table[x=papers, y expr=\coordindex] {\littop};
\end{axis}
\begin{axis}[
    ttaxis,
    at={($(top.east)+(1.9cm,0)$)}, anchor=west,
    width=0.53\linewidth, height=6.4cm,
    ybar, bar width=6.5pt,
    ymin=75, ymax=101,
    ylabel={best reported test accuracy (\%)},
    xtick=data,
    xticklabels from table={\litbest}{dataset},
    xticklabel style={font=\sffamily\scriptsize, color=ttink, rotate=35, anchor=north east},
    xmajorgrids=false,
    enlarge x limits=0.1,
    legend style={at={(0.5,1.02)}, anchor=south, legend columns=2},
    title={TM vs.\ the baselines TM papers report},
    title style={yshift=13pt},
  ]
  \addplot[ttbar] table[x expr=\coordindex, y=tm] {\litbest};
  \addlegendentry{best TM model}
  \addplot[ttbarcpu] table[x expr=\coordindex, y=baseline] {\litbest};
  \addlegendentry{best non-TM baseline}
\end{axis}
\end{tikzpicture}
